\chapter{48}

\section{Moosfluh - Bettmeralp (CH), Montag, 27. August}



Es war kurz nach halb sieben.
Für den Aufstieg bereit traten sie zu dritt vor den Hof. Reich beladen
standen sie direkt im Übergang von der Nacht zum Tag. Louis schaute
hingerissen um sich. Er liebte diese winzig kurze Zeitspanne, bevor sich
die Dunkelheit zugunsten des Lichts ergab. Manuela und Joh waren neben
Louis stehengeblieben. Sie schauten ihm zu. Und warteten.

Nach wenigen Schritten in Richtung Bahnstation \emph{Betten
Dorf} zeigte sich das erste Gelborange der aufgehenden Sonne. Die
filigranen Strahlen erschienen Louis wie vorauseilende Sonnenlaser,
bevor ein dunkles Gold die Landschaft zum Tag erweckte.
Aufgrund des schweren Gepäcks würden sie mit der Gondelbahn nach oben
fahren. Erst auf die Bettmeralp. Und schliesslich nach der kurzen Wanderung zum Blausee mit der Gondel bis hoch zur \emph{Moosfluh}. Selbst von hier aus würden sie
noch gute zwei Stunden gehen müssen, um die Sennhütten zu erreichen.

Oben angekommen halfen sie sich gegenseitig, 
das Gepäck an ihren Körpern festzuschnallen. Eine beeindruckende
Rundumsicht tat sich auf.

«Wahnsinn,» Louis’ Stimme bebte, «diese Stille, diese Schönheit, das
ist, ehm, als wären wir in einer Parallelwelt angekommen.»

Joh legte ihm die Hand auf die Schulter.

«Auch ein Teil unserer Erde, ja,» lächelte er und Louis war sich nicht
sicher, ob er mit sich selbst oder mit ihnen sprach, «die Berge sind der
Inbegriff für die Erhaltung landschaftlicher Ästhetik oder der
ursprünglichen Wirklichkeit. Dieser Unterschied zum Tal,» Joh schüttelte
den Kopf.

Manuela stand unterdessen auf einer flachen Felsplatte. Sie streckte ihr
Gesicht der Sonne entgegen.

«Nun, Ciro, kommt der anstrengende Teil. Denk dran, wir sind keine
beladenen Esel, denen mit dem Stock auf den Arsch gehauen wird,» lachte
Joh und stellte gleich wieder auf Ernsthaftigkeit um, «das wird keine
Übung für Schnelligkeit. Wir haben so viel Gewicht auf uns, dass es um
Konzentration und nicht zuletzt um Willenskraft geht. Wir können
jederzeit Pausen einlegen. Wir steigen nahe beisammen und langsam ab,
alles okay? Ich geh voraus.»

Bereits die ersten Wanderabschnitte hatten Louis Respekt
eingeflösst. Schwer beladen zu gehen, von nun an mehrheitlich zu
klettern, war eine neue Erfahrung. Er nickte ehrfurchtsvoll in Johs
Richtung.
Ihr Weg führte erst über eine kleine Erhöhung. Dann steil und über grosse
Steine hinweg hinunter zum Gletscher.
Jeden einzelnen Schritt musste Louis bewusst und aufmerksam setzen.
Rasch verstand er, was Joh mit geeignetem Schuhwerk gemeint hatte. Der
Bergschuh verlieh ihm Sicherheit. Kein Herumrutschen im Innern des
Schuhs. Jeder Tritt sass. Der bissige Druck auf seinen Rücken dehnte
sich mit jeder Bewegung auf seine Beine aus. Und der Schweiss rann
alsbald in schmalen Rinnsalen über sein Gesicht.
Beim Einatmen war Louis, als stosse der Berg einen betörenden Eigenduft
aus. Die Frische des Lichts war nicht zu überbieten. Und während sich
das Flirren der Seilbahn langsam ausblendete, erinnerten nur mehr ihre knirschenden Bergschuhe bei jedem Schritt daran, dass sie sich hier oben als menschliche Wesen bewegten. 

Louis befand sich an der Schwelle zu neuem Land. Zu ganz und gar
unbekanntem Terrain. Er war auf einen Parallelweg eingebogen, der ihn
mit nichts an sein Leben im Tal erinnerte. Er liess seinen
Blick immer wieder bedächtig vom Boden zum Himmel gleiten.

Schliesslich war der Abstieg geschafft. Joh und Manuela legten ihr
Gepäck ab. Louis’ Knie schlotterten.
Der Himmel hing nun in flächendeckend dicht gespritztem Azurblau über ihnen.
Hier und dort blitzten immer wieder winzig kleine Glitzerpunkte auf.
Flüchtig keck, als spielten sie Verstecken.

«Zauberhaft,» flüsterte Louis vor sich hin und sog die Schönheit des
Ortes in sich ein. 

Er liess seinen Blick der Gletscherzunge entlang von
unten bis weit hinauf gleiten. So etwas hatte er noch nie gesehen. Und
ohne sein Gepäck abzulegen, rückte er direkt an den Rand der Eisschicht.
Dann ging er umständlich in die Knie und schaute auf seine Hände, die
über die kühlen Rundungen strichen.
Joh und Manuela hätte er fast vergessen. Die beiden waren dabei, sich
für eine Pause einzurichten. Louis war fasziniert. Er ging einige
Schritte der Gletscherzunge entlang nach oben und wieder zurück. Als er
sein Gepäck schliesslich neben den beiden abstellte, fielen ihm die
Kopfhörer auf. Joh wirkte abwesend. Was er sich wohl anhörte, hier an
diesem Ort und in diesem Moment?
Manuela lächelte Louis zu. Ihre Augen waren durch die Sonnenbrille nur
als flüchtige Umrisse zu erkennen.

«Der Abstieg lief gut, Ciro, nicht?» sie trat einen Schritt zurück und
setzte die Kaffeetasse an ihre Lippen, «bediene dich.»

Manuela zeigte auf ein buntes Tuch am Boden, bückte sich und griff nach
einem Croissant.

«Ihr habt ein komplettes Frühstück mitgebracht?» freute sich Louis,
«vielen Dank.»

Manuela nickte ihm kauend zu.
Von der Seite konnte Louis erkennen, dass Johs Augen hinter der
Sonnenbrille geschlossen waren. Er schien nach wie vor weit weg zu sein.
Louis goss sich Kaffee aus einer Thermoskanne in die Tasse.

«Mensch, tut das gut,» er hielt die Tasse mit beiden Händen umschlossen,
«ja, der Abstieg war anstrengend, mit all dem Gepäck. Jetzt verstehe
ich, warum Spezialschuhe notwendig sind.»

Manuela biss ein weiteres Stück ihres Hörnchens ab.

«Ich find’s schön, Ciro,» sagte sie, bevor sie weiter ass, «dass du
deinen ersten Alptag unter diesen grossartigen Bedingungen erlebst,» sie
strahlte, als imitiere sie das Licht, «bei diesem Wetter ist die Alp mit
ihrer Umgebung unerreicht, nicht in Worte zu fassen. Manchmal könnte ich
vor Freude heulen.»

«Mir geht’s genauso. Ich kann kaum glauben, dass ich in echt hier stehe,
zum ersten Mal direkt an einem Gletscher.»

«Ach so, um so aufregender für dich, eine Première,» rief sie nun aus
und strich Louis über den Rücken.

«Bald wird auch Joh wieder bei uns sein,» sie setzte sich auf einen
Stein, «kürzlich hat ihm Mazi einen neu erschienen Song empfohlen, den
Joh begeistert zu seiner \emph{Alpmusik} erkor. Und wenn’s nicht aus
Eimern giesst, stellt er sich hier in die Landschaft und hört sich ihn
an. Joh meinte, Mazi habe ihm das Stück genau an diesem Ort das erste
Mal vorgespielt und darum gehöre es hier hin. Der Song hat einen sehr
speziellen Hintergrund, der Joh gefällt. Ort, Musik und Sonne seien
eins, sagt Joh,» sie liess ihre Augen geheimnisvoll kreisen, «ausser
hier, spielt er ihn nie ab.»

Als hätte Joh Manuelas Worte gehört, zog er sich die Kopfhörer aus den
Ohren, setzte ein breites Grinsen auf und streckte seine Arme weit aus.
Manuela stand auf. Er umschlang sie mit einer Freude, die auf Louis
übersprang.

«Ist es nicht paradiesisch hier, Ciro?»

«Unglaublich, ja,» gab Louis zurück und drehte sich zu Joh um, «was hast
du dir gerade angehört? Manuela erzählte mir von deiner
\emph{Alpmusik.»}

«Ach ja? Dacht ich mir schon.»

Joh zog Manuela näher an sich heran. Sie hob ihre Tasse in die Höhe, und wich mit den Worten \emph{Achtung, mein Kaffee} zurück. Joh liess sie schmunzelnd los.

«Ein Geschenk von Mazi sozusagen,» Joh wandte sich Louis zu und hörte
sich hingerissen an, «\emph{Manitoumani»} von «Matthieu Chedid»
und seinen afrikanischen Freunden, eine Art Hommage an die
\emph{Kora}, an die westafrikanische Stegharfe, kennst du das
Instrument?»
  
«Hm, meinst du die Harfenlaute?»

Louis’ Interesse wuchs.

«Ja, kann sein, ich kenne mich mit afrikanischen Instrumenten nicht
besonders aus, auch nicht mit diesem Musikgenre. Mazi ist sehr bewandert
darin und liebt die Art Musik.»

Damit liess Joh sein Handy in der Hosentasche verschwinden. Er ging in
die Knie und griff nach einer Birne.
